% glossary.tex -- Module 8 ESSENTIALS glossary (A-Z): the key terms only.
% src/mod8/glossary.tex (term, meaning and Chinese copied word for word).
% Edit the script, not this file.
\unnumberedchapter[\faBook]{Glossary}

The key words of this module, from A to Z. The meaning is in plain English. The last column
gives the usual Chinese term used in physics textbooks. The full guide's glossary has more
words.

\begin{glossarytable}
  % ---- A
  \gterm{Absolute magnitude ($M$)}{brightness at \qty{10}{pc}; smaller = brighter}{绝对星等}
  \gterm{Absorption spectrum}{dark lines on a continuous spectrum}{吸收光谱}
  \gterm{Activity ($A$)}{decays per second, in becquerel (Bq)}{放射性活度}
  \gterm{Alpha particle ($\alpha$)}{a helium-4 nucleus, charge $+2e$}{α粒子}
  \gterm{Annihilation}{particle + antiparticle turn into photons}{湮灭}
  \gterm{Antiparticle}{same mass; opposite charge and properties}{反粒子}
  \gterm{Artificial transmutation}{a nuclear change forced by firing a particle}{人工核转变}
  \gterm{Atomic number ($Z$)}{number of protons; it decides the element}{原子序数}
  % ---- B
  \gterm{Balmer series}{hydrogen lines ending on $n = 2$ (four visible)}{巴耳末系}
  \gterm{Band of stability}{the strip of stable nuclei on an $N$--$Z$ graph}{稳定带}
  \gterm{Baryon}{a hadron of three quarks, such as the proton}{重子}
  \gterm{Beta particle ($\beta$)}{a fast electron ($\beta^-$) or positron ($\beta^+$)}{β粒子}
  \gterm{Big Bang}{the hot, dense start of the Universe}{大爆炸}
  \gterm{Big Bang nucleosynthesis}{helium nuclei made 3--20 minutes after it}{大爆炸核合成}
  \gterm{Binding energy ($E_\mathrm{B}$)}{energy needed to pull a nucleus apart}{结合能}
  \gterm{Binding energy per nucleon}{$E_\mathrm{B} \div A$; bigger = more tightly bound}{比结合能}
  \gterm{Black body}{ideal absorber; its glow depends only on $T$}{黑体}
  \gterm{Boson}{a particle with whole-number spin}{玻色子}
  % ---- C
  \gterm{Cathode ray}{a stream of electrons from the cathode}{阴极射线}
  \gterm{Cepheid variable}{a pulsing star; its period gives its luminosity}{造父变星}
  \gterm{Chain reaction}{neutrons from fissions cause more fissions}{链式反应}
  \gterm{Charge-to-mass ratio ($q/m$)}{charge divided by mass, in \si{C.kg^{-1}}}{比荷}
  \gterm{CNO cycle}{fusion of H to He with C, N, O as catalysts}{碳氮氧循环}
  \gterm{Continuous spectrum}{every wavelength, with no gaps}{连续光谱}
  \gterm{Control rod}{a rod (boron, cadmium) that absorbs neutrons}{控制棒}
  \gterm{Cosmic microwave background (CMB)}{faint microwaves from every direction}{宇宙微波背景辐射}
  \gterm{Critical mass}{least mass that keeps a chain reaction going}{临界质量}
  \gterm{Cyclotron}{particles spiral out in a fixed magnetic field}{回旋加速器}
  % ---- D
  \gterm{de Broglie wavelength}{the wavelength of a moving particle, $\frac{h}{mv}$}{德布罗意波长}
  \gterm{Decay constant ($\lambda$)}{chance per second that one nucleus decays}{衰变常量}
  % ---- E
  \gterm{Elementary charge ($e$)}{smallest free charge, \qty{1.602e-19}{C}}{元电荷}
  \gterm{Emission spectrum}{bright lines at a few exact wavelengths}{发射光谱}
  \gterm{Energy level}{an allowed energy of an electron in an atom}{能级}
  \gterm{Exchange particle}{a boson that carries a force}{交换粒子}
  % ---- F
  \gterm{Fermion}{a particle with half-number spin}{费米子}
  \gterm{Fundamental particle}{a particle with no smaller parts}{基本粒子}
  % ---- G
  \gterm{Gamma ray ($\gamma$)}{a high-energy photon from a nucleus}{γ射线}
  \gterm{Ground state}{the lowest energy level ($n = 1$)}{基态}
  % ---- H
  \gterm{Hadron}{a particle made of quarks}{强子}
  \gterm{Half-life ($t_{\frac{1}{2}}$)}{time for half of the nuclei to decay}{半衰期}
  \gterm{Hertzsprung--Russell (HR) diagram}{stars' luminosity against temperature}{赫罗图}
  \gterm{Higgs boson}{a ripple in the Higgs field, found in 2012}{希格斯玻色子}
  \gterm{Hubble's law}{recession speed $\propto$ distance: $v = H_0 d$}{哈勃定律}
  % ---- I
  \gterm{Ionising radiation}{radiation that knocks electrons out of atoms}{电离辐射}
  \gterm{Isotope}{same $Z$, different number of neutrons}{同位素}
  % ---- L
  \gterm{Lepton}{a fundamental particle with no strong force}{轻子}
  \gterm{Luminosity ($L$)}{total power a star gives out, in watts}{光度}
  \gterm{Lyman series}{hydrogen lines ending on $n = 1$ (UV)}{莱曼系}
  % ---- M
  \gterm{Main sequence}{the long stage when a star fuses H into He}{主序}
  \gterm{Mass defect ($\Delta m$)}{mass lost as nucleons join, or in a reaction}{质量亏损}
  \gterm{Mass number ($A$)}{number of protons plus neutrons}{质量数}
  \gterm{Meson}{a hadron of one quark and one antiquark}{介子}
  \gterm{Moderator}{material that slows neutrons down}{慢化剂}
  % ---- N
  \gterm{Neutrino}{a neutral lepton with a tiny mass}{中微子}
  \gterm{Neutron}{a neutral particle in the nucleus}{中子}
  \gterm{Nuclear fission}{a heavy nucleus splits into two smaller ones}{核裂变}
  \gterm{Nuclear fusion}{light nuclei join to make a heavier one}{核聚变}
  \gterm{Nuclear model}{Rutherford's atom, with a tiny, dense nucleus}{核式结构模型}
  \gterm{Nucleon}{a proton or a neutron}{核子}
  % ---- O
  \gterm{Orbital}{a region where an electron is likely to be}{原子轨道}
  % ---- P
  \gterm{Pair production}{a photon becomes a particle--antiparticle pair}{对产生}
  \gterm{Paschen series}{hydrogen lines ending on $n = 3$ (IR)}{帕邢系}
  \gterm{Photon}{a packet of light energy, $E = hf$}{光子}
  \gterm{Plum pudding model}{Thomson's atom: electrons in a $+$ ball}{枣糕模型}
  \gterm{Positron}{the electron's antiparticle, charge $+e$}{正电子}
  \gterm{Proton--proton chain}{fusion of 4 protons into He-4 (Sun-like stars)}{质子-质子链}
  % ---- Q
  \gterm{Quantised}{only certain separate values are allowed}{量子化的}
  \gterm{Quark}{a fundamental particle inside hadrons}{夸克}
  % ---- R
  \gterm{Radioactive decay}{an unstable nucleus changes by itself}{放射性衰变}
  \gterm{Red giant}{a big, cool, bright star late in its life}{红巨星}
  \gterm{Redshift}{lines moved to longer wavelengths}{红移}
  % ---- S
  \gterm{Spectral class}{O B A F G K M, from hot to cool}{光谱型}
  \gterm{Standard Model}{theory of fundamental particles and forces}{标准模型}
  \gterm{Stationary state}{an allowed orbit, with no radiation}{定态}
  \gterm{Strong force}{holds quarks together; carried by gluons}{强相互作用}
  \gterm{Strong nuclear force}{short-range attraction between nucleons}{核力}
  % ---- T
  \gterm{Transmutation}{a nuclear change to another element}{核嬗变}
  % ---- U
  \gterm{Uncertainty principle}{position and momentum are not both exact}{不确定性关系}
  % ---- W
  \gterm{W and Z bosons}{the exchange particles of the weak force}{W和Z玻色子}
  \gterm{Wave function ($\psi$)}{the electron's wave; $\psi^2$ gives the probability}{波函数}
  \gterm{Weak force}{changes quark flavour, as in beta decay}{弱相互作用}
  \gterm{White dwarf}{small, hot, dim core left by a Sun-like star}{白矮星}
  \gterm{Wien's displacement law}{$\lambda_\mathrm{max} = \frac{b}{T}$: hotter means shorter}{维恩位移定律}
\end{glossarytable}
